% Converted from sections/00-abstract.md by tools/md2tex.py; edit this file from now on.
Declarative infrastructure-as-code (IaC) tools reconcile a desired state with the observed one,
which presumes APIs that expose resources with stable identities, standard methods and typed
representations. Many self-hosted platforms expose procedures instead. Dokploy, an open-source
deployment platform, offers 604 procedures through a tRPC API whose OpenAPI export declares every
response as an empty object; it writes one object through several procedures and often does not
return the identity of what it creates. This report describes a method for building a declarative
control plane over such an API without changing it. What the API leaves implicit is written once per
resource kind, as data, in a kind specification: identity, the role of each procedure, the fields
each update procedure carries, how a new identity is learned, which reads prove absence, the
mutability and secrecy of each field, and shapes such as unions, relations and kinds with several
parents. A kind-agnostic engine derives parsing, discovery, three-way planning, journaled execution
with read-back verification, recovery and import from these specifications. A coverage ledger checks
completeness by classifying every request field of every procedure used against the API contract,
and response shapes are taken from recorded live responses. For Dokploy, 36 specifications with 625
fields use 223 of the 604 procedures, and no request field is left unclassified. \result{convergence, idempotency, recovery and round-trip outcomes.} The report also identifies twelve
classes of mismatch between procedure-oriented APIs and declarative management and gives a procedure
for applying the method to other APIs.

\textbf{Keywords:} infrastructure as code, declarative configuration, reconciliation, remote procedure call,
OpenAPI, specification-driven engineering, Dokploy
